\documentclass[blue]{beamer}
\mode<presentation> {
  \setbeamertemplate{background canvas}[vertical shading][bottom=red!10,top=blue!10]
  \usetheme{Boadilla}
  \usefonttheme[onlysmall]{serif}
}
\usepackage{pgf,pgfarrows,pgfnodes,pgfautomata,pgfheaps,pgfshade}
\usepackage{amsmath,amssymb,bm}
\usepackage{colortbl}
\usepackage[english]{babel}
\setbeamercovered{dynamic}
\def\R{\mathbb{R}}
\def\Q{\mathbb{Q}}
\def\N{\mathbb{N}}
\def\E{\mathbb{E}}
\newcommand{\ps}[2]{ \langle #1 , #2 \rangle }
\def\1{\mathbf{1}}
\newcommand{\nor}[1]{\left\| #1 \right\|}
\def\leq{\leqslant}
\def\geq{\geqslant}
\def\et{\quad \textrm{and} \quad}
\def\ou{\quad \textrm{or} \quad}
\def\boite{\hskip 0.5mm {\Box} \hskip 0.5mm}
\DeclareMathOperator\PPE{PPE}
\DeclareMathOperator\vNM{vNM}
\DeclareMathOperator\argmax{argmax}
\DeclareMathOperator\Prob{Prob}
\DeclareMathOperator\supp{supp}
\DeclareMathOperator\Eq{Eq}
\DeclareMathOperator\inte{int}
\DeclareMathOperator\co{co}
\DeclareMathOperator\cl{cl}
\DeclareMathOperator\Dirac{Dirac}
\DeclareMathOperator\F{F}
\DeclareMathOperator\As{As}
\DeclareMathOperator\Ir{Ir}
\DeclareMathOperator\MRS{MRS}
\DeclareMathOperator\Range{Im}
\DeclareMathOperator\DIV{DIV}
%------------------------------ theorem styles
\theoremstyle{definition}
\newtheorem*{assumption}{Assumption}
\newtheorem*{conjecture}{Conjecture}
\theoremstyle{example}
\newtheorem*{proposition}{Proposition}
\newtheorem*{theoremns}{Theorem}
\newtheorem*{remark}{Remark}
\newtheorem*{definitionns}{Definition}
\newtheorem*{lemmans}{Lemma}
%-----------------------------------------
\title[Corporate Finance]{Theory of Corporate Finance}
\author[Victor Filipe]{V. Filipe Martins-da-Rocha}
\institute[FGV EESP]{Escola de Economia de S\~ao Paulo}
\date[April--June, 2026]{Part 1.d. Stock Market Equilibrium: Moral hazard}
\subject{Economic Theory}
\begin{document}
\frame{\titlepage}

%--------------------------------------------------------------------
\begin{frame}\frametitle{Outline}
\begin{itemize}
\item The framework: entrepreneurs, investors, securities
\item Optimal effort and inference
\item Stock-market equilibrium with price perceptions
\item Rational competitive price perceptions (RCPP)
\item Partial spanning: definition and examples
\item Constrained Pareto efficiency
\item Strong partial spanning and the welfare theorem
\item Limitations
\end{itemize}
\end{frame}

%====================================================================
\section{The Question}
%====================================================================

\begin{frame}\frametitle{The question at stake}
\begin{itemize}
\item Stock markets allocate resources (constrained) efficiently when there are no frictions
\item Is this still true in the presence of moral hazard?
\bigskip
\item Consider an entrepreneur whose effort affects the firm's output
\item Effort is not contractible (e.g., not observable)
\end{itemize}
\bigskip

\textbf{The Jensen--Meckling tension (1976):}
\begin{itemize}
\item Optimal risk-sharing may require the entrepreneur to sell part of his firm
\item This may reduce his incentives to exert effort
\end{itemize}
\bigskip

First-best is generally not achievable; the right question is about \emph{constrained} efficiency
\end{frame}

%====================================================================
\section{The Framework}
%====================================================================

\begin{frame}\frametitle{Time, agents, and goods}
\begin{itemize}
\item Two-period economy with production, single commodity per period
\item Uncertainty at $t=1$: a finite set $S$ of states of nature
\end{itemize}
\bigskip

The set of agents is partitioned:
\begin{equation*}
I = I_1 \cup I_2
\end{equation*}

\begin{itemize}
\item $I_1$: finite set of \emph{entrepreneurs}
\item $I_2$: finite set of \emph{investors}
\end{itemize}
\bigskip

Each agent has initial wealth $\omega^i_0 \geq 0$ at $t=0$ and zero endowment at $t=1$
\end{frame}

%--------------------------------------------------------------------
\begin{frame}\frametitle{The production technology}
For each entrepreneur $i \in I_1$:
\begin{itemize}
\item invest capital $k^i \geq 0$ at $t=0$
\item exert effort $e^i$ in an abstract set $E^i$
\item produce in state $s$:
\begin{equation*}
f^i_s(k^i, e^i) \geq 0
\end{equation*}
\end{itemize}
\bigskip

We write $f^i(k^i, e^i) = (f^i_s(k^i, e^i))_{s \in S}$
\bigskip

For investors $i \in I_2$: $E^i = \{0\}$ and $f^i \equiv 0$
\end{frame}

%--------------------------------------------------------------------
\begin{frame}\frametitle{The financing question}
Entrepreneurs need funds to finance investment
\bigskip

Since production is uncertain, risk-sharing between entrepreneurs and investors is desirable
\bigskip

Investors transfer wealth by buying stocks of firms
\bigskip

\textbf{Tension:} If the entrepreneur sells too much equity, his wealth no longer depends on his effort, and he won't exert effort
\bigskip

\textbf{The new shareholders} of this firm will then receive a low dividend
\bigskip

\begin{block}{Question}
Under what conditions can markets satisfactorily solve the joint problem of financing, risk-sharing, and incentives?
\end{block}
\end{frame}

%--------------------------------------------------------------------
\begin{frame}\frametitle{Available financial markets}
\textbf{What financial markets exist in a modern economy?}
\bigskip

\begin{assumption}[Effort is not contractible]
There is no explicit market on which entrepreneurial effort is bought or sold
\end{assumption}
\bigskip

\textbf{Why?}
\begin{itemize}
\item Effort is a complex input, not readily observable or measurable
\item Contracts contingent on effort are difficult to enforce
\end{itemize}
\end{frame}

%--------------------------------------------------------------------
\begin{frame}\frametitle{Available financial markets (cont.)}
\begin{assumption}[No state-contingent contracts]
There are no markets for contracts contingent on the realization of the primitive states of nature
\end{assumption}
\bigskip

\textbf{Why?}
\begin{itemize}
\item Some standard insurable risks exist (fire in a warehouse), but we focus on \emph{business risks} (technological shocks)
\item It is hard to describe ex-ante and verify ex-post the precise nature of these primitive shocks
\item Transaction costs of itemizing contingencies make trading state-contingent contracts unfeasible
\end{itemize}
\end{frame}

%--------------------------------------------------------------------
\begin{frame}\frametitle{The financial contracts}
The contracts available for financing firms and sharing risks:
\bigskip

\begin{itemize}
\item \textbf{Bonds:} non-contingent contracts
\item \textbf{Output-based contracts:} contingent on the (observable) realized outputs of firms
\end{itemize}
\bigskip

\begin{block}{Examples}
\begin{itemize}
\item equity of firms (stocks)
\item derivative securities (e.g., options on equity)
\end{itemize}
\end{block}
\end{frame}

%--------------------------------------------------------------------
\begin{frame}\frametitle{Securities: setup}
There is a finite set $J$ of securities. Each $j \in J$ has a payoff function
\begin{equation*}
V_j : \R^I \to \R
\end{equation*}

Let $y = (y^i)_{i \in I}$ denote the random outputs of firms (with $y^i = 0$ for $i \in I_2$)
\bigskip

\begin{itemize}
\item $y_s = (y^i_s)_{i \in I}$: cross-section of outputs in state $s$
\item $V_j(y_s)$: payoff of security $j$ in state $s$
\item $V_j(y) = (V_j(y_s))_{s \in S}$: state-contingent payoff vector
\end{itemize}
\end{frame}

%--------------------------------------------------------------------
\begin{frame}\frametitle{Marketed subspace}
The linear mapping $V(y) : \R^J \to \R^S$ is defined by
\begin{equation*}
V(y) z := \sum_{j \in J} z_j V_j(y), \quad z \in \R^J
\end{equation*}

The \emph{marketed subspace} at output $y$ is
\begin{equation*}
\mathcal{V}(y) := \Range V(y) \subset \R^S
\end{equation*}

This is the set of state-contingent claims at $t=1$ that can be replicated by a portfolio of traded securities
\end{frame}

%--------------------------------------------------------------------
\begin{frame}\frametitle{Four classes of securities}
\textbf{Bond.} One asset $j_0$ with $V_{j_0}(y) = \1_S$ for every $y$
\bigskip

\textbf{Equity.} For each firm $i \in I_1$, a security with $V_i(y) = y^i$
\bigskip

\textbf{Simple options.} A call on firm $i$'s equity with strike $\tau$:
\begin{equation*}
V_j(y_s) = \max\{ y^i_s - \tau, 0 \}
\end{equation*}


\textbf{Indices and complex options.} An index is a weighted sum:
\begin{equation*}
V_j(y_s) = \sum_{i \in I} \gamma^i y^i_s, \quad \sum_i \gamma^i = 1
\end{equation*}
A complex option is a call/put on an index
\end{frame}

%--------------------------------------------------------------------
\begin{frame}\frametitle{Inside vs.\ outside securities}
For each entrepreneur $i \in I_1$, partition the securities:
\begin{equation*}
J = J^i \cup J^i_{-i} \cup J_{-i}
\end{equation*}

\begin{itemize}
\item $J^i$: securities depending \emph{exclusively} on firm $i$'s output
    \begin{itemize}
    \item equity $i$, simple options on $i$
    \end{itemize}
\item $J^i_{-i}$: securities depending on $i$'s output \emph{and} on at least one other firm
    \begin{itemize}
    \item indices, complex options
    \end{itemize}
\item $J_{-i}$: securities not depending on firm $i$'s output
    \begin{itemize}
    \item bond, equity of other firms, options on other firms
    \end{itemize}
\end{itemize}
\end{frame}

%====================================================================
\section{Budgets and Effort}
%====================================================================

\begin{frame}\frametitle{Entrepreneur's budget}
Given $q$, $y^{-i}$, and choices $(k^i, z^i, e^i)$:
\bigskip

\begin{block}{At $t = 0$}
\begin{equation*}
0 \leq c^i_0 := \omega^i_0 + q_i - q \cdot z^i - k^i
\end{equation*}
\end{block}

\begin{block}{At $t = 1$ in state $s$}
\begin{equation*}
0 \leq c^i_s := V_s(f^i(k^i, e^i), y^{-i}) \cdot z^i
\end{equation*}
\end{block}
\bigskip

The entrepreneur receives $q_i$ from selling his equity, uses the proceeds to finance investment $k^i$, and consumes the residual
\end{frame}

%--------------------------------------------------------------------
\begin{frame}\frametitle{Investor's budget}
With $k^i = 0$, $e^i = 0$, $q_i = 0$, $f^i = 0$:
\bigskip

\begin{block}{At $t = 0$}
\begin{equation*}
0 \leq c^i_0 := \omega^i_0 - q \cdot z^i
\end{equation*}
\end{block}

\begin{block}{At $t = 1$ in state $s$}
\begin{equation*}
0 \leq c^i_s := V_s(0, y^{-i}) \cdot z^i
\end{equation*}
\end{block}
\bigskip

\textbf{No short-sale constraints in this chapter:} default penalties are assumed to be so harsh that all personal debts are honoured
\end{frame}

%--------------------------------------------------------------------
\begin{frame}\frametitle{Preferences}
Each agent $i$ has utility $U^i : \R_+ \times \R_+^S \times E^i \to \R$
\bigskip

\begin{assumption}[Time-separability]
\begin{equation*}
U^i(c^i_0, c^i_1, e^i) = u^i_0(c^i_0) + u^i_1(c^i_1, e^i)
\end{equation*}
\end{assumption}
\bigskip

with $u^i_0$ continuous, strictly increasing, strictly concave; $u^i_1$ continuous
\bigskip

Time-separability ensures that the entrepreneur's incentive to exert effort is driven exclusively by the effect of effort on his $t=1$ consumption
\end{frame}

%--------------------------------------------------------------------
\begin{frame}\frametitle{Timing of decisions}
\textbf{Crucial timing assumption:}
\bigskip

The entrepreneur chooses effort $e^i$ \alert{after} the financial decisions $(k^i, z^i)$ have been made, but \alert{before} the state of nature is realized
\bigskip

\begin{itemize}
\item He is committed to $(k^i, z^i)$ when he picks $e^i$
\item He cannot re-trade securities after observing effort-induced expected payoffs
\end{itemize}
\bigskip

The capital structure of the firm thus affects the performance of management through incentives
\end{frame}

%--------------------------------------------------------------------
\begin{frame}\frametitle{The optimal effort correspondence}
For given $(k^i, z^i, y^{-i})$, define
\begin{equation*}
\widetilde{e}^i(k^i, z^i, y^{-i}) := \argmax \bigl\{ u^i_1(c^i_1, e^i) \colon e^i \in E^i,\ c^i_1 = V(f^i(k^i, e^i), y^{-i}) z^i \bigr\}
\end{equation*}

\begin{assumption}[Existence of optimal effort]
The correspondence $\widetilde{e}^i(\cdot)$ has non-empty values for every $(k^i, z^i, y^{-i})$.
\end{assumption}
\bigskip

Every agent agrees that entrepreneur $i$ will choose some $e^i \in \widetilde{e}^i(k^i, z^i, y^{-i})$
\end{frame}

%====================================================================
\section{Inference and Price Perceptions}
%====================================================================

\begin{frame}\frametitle{Anticipating efforts}
Consider an investor or entrepreneur who is thinking of buying equity or options of firm $i$
\bigskip

\textbf{To anticipate the firm's profit,} the agent needs to anticipate $(k^i, e^i)$
\bigskip

\textbf{The capital input $k^i$ is observable.}
\bigskip

\textbf{The effort $e^i$ is unobservable,} but it can be \emph{deduced} if the entrepreneur's characteristics $(u^i_1, f^i)$ and his financial variables are known
\end{frame}

%--------------------------------------------------------------------
\begin{frame}\frametitle{Information assumption}
\begin{assumption}[Common knowledge of characteristics]
For every entrepreneur $i \in I_1$, the tuple $(u^i_1, f^i, k^i, z^i)$ is common knowledge among all agents.
\end{assumption}
\bigskip

\textbf{Plausibility (Magill--Quinzii 2002):}
\begin{itemize}
\item SEC disclosure rules require statements on capital projects, equity holdings, options
\item Securities analysts accumulate detailed knowledge of firms and managers
\item Brokerage firms disseminate this information
\end{itemize}
\bigskip

The informational asymmetry at the level of effort is compensated by full observability of financial decisions
\end{frame}

%--------------------------------------------------------------------
\begin{frame}\frametitle{Entrepreneurs as signallers}
Investors infer effort from $(k^i, z^i)$ and price securities accordingly
\bigskip

Entrepreneurs understand this and use their financial choices as \emph{signals}:
\begin{itemize}
\item retaining a high equity stake signals high effort
\item full diversification signals low effort
\end{itemize}
\bigskip

\textbf{In equilibrium:} the signal is correct
\begin{itemize}
\item investors' inference matches the entrepreneur's actual effort
\item the entrepreneur, anticipating this, makes the financial choice that best trades off incentives against risk-sharing
\end{itemize}
\end{frame}

%--------------------------------------------------------------------
\begin{frame}\frametitle{Price perceptions: notation}
For each agent $i \in I$ and each security $j \in J$:
\begin{equation*}
\widetilde{Q}^i_j(k^i, z^i; \overline{y}^{-i}) \in \R
\end{equation*}
\bigskip

denotes the price agent $i$ expects for security $j$, conditional on:
\begin{itemize}
\item entrepreneur $i$'s alternative financial choice $(k^i, z^i)$
\item the equilibrium outputs $\overline{y}^{-i}$ of all other firms
\end{itemize}
\bigskip

The collection $\widetilde{Q}^i = (\widetilde{Q}^i_j)_{j \in J}$ is agent $i$'s price perception
\end{frame}

%====================================================================
\section{Equilibrium with Price Perceptions}
%====================================================================

\begin{frame}\frametitle{Equilibrium: definition}
\begin{definitionns}
A stock-market equilibrium with price perceptions $\widetilde{Q}$ is a list
\begin{equation*}
\bigl\{ (\overline{c}^i, \overline{e}^i, \overline{k}^i, \overline{z}^i)_{i \in I}, (\overline{y}^j)_{j \in J}, \overline{q} \bigr\}
\end{equation*}
satisfying conditions (a)--(d) on the next slides.
\end{definitionns}
\end{frame}

%--------------------------------------------------------------------
\begin{frame}\frametitle{Equilibrium: condition (a) -- agents are rational}
For each $i \in I$, $(\overline{c}^i, \overline{e}^i, \overline{k}^i, \overline{z}^i)$ maximizes $U^i(c^i_0, c^i_1, e^i)$ subject to:
\bigskip

\begin{equation*}
c^i_0 = \omega^i_0 + \widetilde{Q}^i_i(k^i, z^i; \overline{y}^{-i}) - \widetilde{Q}^i(k^i, z^i; \overline{y}^{-i}) \cdot z^i - k^i
\end{equation*}
\bigskip

\begin{equation*}
c^i_1 = V(f^i(k^i, e^i), \overline{y}^{-i}) z^i
\end{equation*}
\bigskip

The budget set is denoted $B^i(\overline{q}, \overline{y}^{-i})$
\end{frame}

%--------------------------------------------------------------------
\begin{frame}\frametitle{Equilibrium: conditions (b)--(d)}
\begin{block}{(b) Production is compatible with effort}
For every firm $i \in I_1$:
\begin{equation*}
\overline{y}^i = f^i(\overline{k}^i, \overline{e}^i)
\end{equation*}
\end{block}

\begin{block}{(c) Price perceptions are consistent at equilibrium}
For every agent $i \in I$:
\begin{equation*}
\overline{q} = \widetilde{Q}^i(\overline{k}^i, \overline{z}^i; \overline{y}^{-i})
\end{equation*}
\end{block}

\begin{block}{(d) Markets clear}
\begin{equation*}
\sum_{i \in I} \overline{z}^i_j = 1 \text{ for $j \in I$}, \quad \sum_{i \in I} \overline{z}^i_j = 0 \text{ for $j \notin I$}
\end{equation*}
\end{block}
\end{frame}

%--------------------------------------------------------------------
\begin{frame}\frametitle{What the equilibrium captures}
\begin{itemize}
\item each entrepreneur takes the production plans and security prices of other firms as given
\item each entrepreneur anticipates how his financial decisions influence the prices of securities depending on his firm's outputs
\item the output anticipated by outside investors is compatible with the entrepreneur's choice of effort
\item the price perceptions are consistent with the observed equilibrium prices
\end{itemize}
\bigskip

\begin{block}{Open question}
How to define $\widetilde{Q}^i(k^i, z^i; \overline{y}^{-i})$ \emph{outside} the equilibrium path?
\end{block}
\end{frame}

%====================================================================
\section{Pricing Out-of-Equilibrium}
%====================================================================

\begin{frame}\frametitle{What entrepreneur $i$ needs to predict}
To form $\widetilde{Q}^i(k^i, z^i; \overline{y}^{-i})$, entrepreneur $i$ needs:
\bigskip

\begin{enumerate}
\item[(a)] the \alert{output} of his firm that investors expect upon observing $(k^i, z^i)$
\bigskip
\item[(b)] how the market will \alert{price} the resulting securities
\end{enumerate}
\bigskip

\textbf{Answer to (a):} Investors deduce from $(k^i, z^i)$ what the likely effort
\begin{equation*}
\widehat{e}^i \in \widetilde{e}^i(k^i, z^i, \overline{y}^{-i})
\end{equation*}
will be, and hence the likely output $f^i(k^i, \widehat{e}^i)$
\end{frame}

%--------------------------------------------------------------------
\begin{frame}\frametitle{State prices at equilibrium}
At equilibrium, no arbitrage implies the existence of state prices $\overline{\pi} \in \R_{++}^S$ such that
\begin{equation*}
\overline{q}_j = \sum_{s \in S} \overline{\pi}_s V_j(\overline{y}_s), \quad \text{for all $j \in J$}
\end{equation*}
\bigskip

\begin{itemize}
\item If markets are complete, $\overline{\pi}$ is unique
\item Otherwise, infinitely many state prices are compatible with $(\overline{q}, \overline{y})$
\end{itemize}
\end{frame}

%--------------------------------------------------------------------
\begin{frame}\frametitle{Pricing of marketed cash flows}
Let $\mathcal{V}(\overline{y})$ be the marketed subspace at equilibrium
\bigskip

For any $m \in \mathcal{V}(\overline{y})$, the value
\begin{equation*}
\overline{\pi} \cdot m = \sum_{s \in S} \overline{\pi}_s m_s
\end{equation*}
is \emph{independent} of the choice of $\overline{\pi}$
\bigskip

We denote this common value by $v(m; \overline{q}, \overline{y})$
\bigskip

\textbf{Why?} Two state prices that disagree on $\R^S$ must agree on $\mathcal{V}(\overline{y})$, since their difference lies in the orthogonal complement of the rows of $V(\overline{y})$
\end{frame}

%--------------------------------------------------------------------
\begin{frame}\frametitle{Pricing a marketed innovation}
Suppose entrepreneur $i$ deviates to $(k^i, z^i)$ and the resulting cash flow $f^i(k^i, \widehat{e}^i)$ \emph{lies in the marketed subspace} $\mathcal{V}(\overline{y})$
\bigskip

Then the entrepreneur may price the new asset using any state-price vector compatible with equilibrium:
\begin{equation*}
\widetilde{Q}^i(k^i, z^i; \overline{y}^{-i}) := \sum_{s \in S} \overline{\pi}_s V_s(f^i(k^i, \widehat{e}^i), \overline{y}^{-i})
\end{equation*}
\bigskip

The result is independent of the choice of $\overline{\pi}$
\bigskip

\textbf{Caveat:} this requires the deviation to stay within $\mathcal{V}(\overline{y})$
\end{frame}

%====================================================================
\section{Partial Spanning}
%====================================================================

\begin{frame}\frametitle{Partial spanning: definition}
\begin{definitionns}[Partial spanning]
There is \alert{partial spanning} (PS) at $\overline{y}$ if for every $i \in I$, every alternative $(k^i, e^i)$ with $y^i = f^i(k^i, e^i)$:
\begin{equation*}
\mathcal{V}(y) \subset \mathcal{V}(\overline{y})
\end{equation*}
\end{definitionns}
\bigskip

\begin{itemize}
\item Under PS, a firm cannot create a ``new security'' by modifying its financial decisions
\item Market prices of existing securities are sufficient signals to value any alternative production plan
\end{itemize}
\end{frame}

%--------------------------------------------------------------------
\begin{frame}\frametitle{RCPP equilibrium}
\begin{definitionns}[RCPP]
A stock-market equilibrium with price perceptions $\widetilde{Q}$ is a \alert{rational competitive price perceptions} (RCPP) equilibrium if:
\end{definitionns}

\begin{itemize}
\item[(i)] partial spanning holds at $\overline{y}$
\bigskip
\item[(ii)] for every $i$ and every alternative $(k^i, z^i)$:
\begin{equation*}
\widetilde{Q}^i_j(k^i, z^i; \overline{y}^{-i}) = \sum_{s \in S} \overline{\pi}_s V_j(f^i(k^i, \widehat{e}^i), \overline{y}^{-i})
\end{equation*}
for any $\overline{\pi}$ compatible with equilibrium and $\widehat{e}^i$ in $\widetilde{e}^i(k^i, z^i, \overline{y}^{-i})$ that maximizes the entrepreneur's date-$0$ consumption
\end{itemize}
\end{frame}

%--------------------------------------------------------------------
\begin{frame}\frametitle{RCPP: rational}
\textbf{Why ``rational''?}
\bigskip

Investors correctly deduce $\widehat{e}^i$ from $(k^i, z^i)$ and the common knowledge of $(u^i_1, f^i)$
\bigskip

\begin{itemize}
\item if $\widetilde{e}^i(k^i, z^i, \overline{y}^{-i})$ is single-valued, the inference is immediate
\item if multi-valued, agents select the element yielding the highest date-$0$ income for the entrepreneur
\end{itemize}
\bigskip

The entrepreneur has every incentive to signal the highest-valuation effort level
\end{frame}

%--------------------------------------------------------------------
\begin{frame}\frametitle{RCPP: competitive}
\textbf{Why ``competitive''?}
\bigskip

Entrepreneurs take state prices as given:
\begin{itemize}
\item an individual firm's deviation does not perturb the equilibrium state-price vector
\item even when the deviation is large relative to the firm's own output
\end{itemize}
\bigskip

This is the price-taking hypothesis adapted to the moral-hazard framework
\end{frame}

%====================================================================
\section{Sufficient Conditions for PS}
%====================================================================

\begin{frame}\frametitle{Single-factor model}
\begin{proposition}[Bond, equity, single-factor risk]
Assume:
\begin{itemize}
\item[(a)] $J = I \cup \{j_0\}$ (only equities and a riskless bond)
\item[(b)] for every $i \in I_1$:
\begin{equation*}
f^i(k^i, e^i) = f^i_c(k^i, e^i) \1_S + g^i(k^i, e^i) \bm{\eta}^i
\end{equation*}
with $f^i_c, g^i$ concave non-decreasing and $\bm{\eta}^i \in \R_+^S$ a fixed risk-structure vector
\end{itemize}
If, at equilibrium, $g^i(\overline{k}^i, \overline{e}^i) > 0$ for every $i \in I_1$, then partial spanning holds at $\overline{y}$.
\end{proposition}
\end{frame}

%--------------------------------------------------------------------
\begin{frame}\frametitle{Single-factor model: idea}
Firm $i$'s output decomposes:
\begin{itemize}
\item deterministic component $f^i_c(k^i, e^i) \1_S$, spanned by the bond
\item risky component $g^i(k^i, e^i) \bm{\eta}^i$, proportional to a fixed direction $\bm{\eta}^i$
\end{itemize}
\bigskip

At equilibrium, the marketed subspace contains:
\begin{itemize}
\item $\1_S$ (the bond)
\item $\bm{\eta}^i = [\overline{y}^i - \overline{f}^i_c \1_S] / \overline{g}^i$ (from the equity)
\end{itemize}
\bigskip

Any alternative output is a linear combination of these spanned vectors, hence in $\mathcal{V}(\overline{y})$
\end{frame}

%--------------------------------------------------------------------
\begin{frame}\frametitle{Independent shocks with options}
\begin{proposition}[Independent shocks plus options]
Assume:
\begin{itemize}
\item[(a)] markets consist of bond, equity, and call/put options on each firm at a rich set of strike prices
\item[(b)] $S = \prod_{i \in I_1} S^i$ and there exist $g^i_{s^i}$ such that
\begin{equation*}
f^i_s(k^i, e^i) = g^i_{s^i}(k^i, e^i)
\end{equation*}
\end{itemize}
If $s^i \mapsto g^i_{s^i}(\overline{k}^i, \overline{e}^i)$ is injective and strikes lie between consecutive output values, then partial spanning holds at $\overline{y}$.
\end{proposition}
\end{frame}

%--------------------------------------------------------------------
\begin{frame}\frametitle{A counter-example: setup}
Consider one entrepreneur and three states: $S = \{s_1, s_2, s_3\}$
\bigskip

Effort set: $E^i = \{e^i_h, e^i_\ell\}$, with $y^i_h > y^i_\ell > 0$
\bigskip

Production does not depend on capital:
\begin{equation*}
f^i(e^i_h) = (y^i_h, y^i_h, y^i_\ell)
\end{equation*}
\begin{equation*}
f^i(e^i_\ell) = (y^i_h, y^i_\ell, y^i_\ell)
\end{equation*}
\bigskip

Securities: bond, equity, simple options
\end{frame}

%--------------------------------------------------------------------
\begin{frame}\frametitle{A counter-example: failure of PS}
At equilibrium with high effort $\overline{e}^i = e^i_h$:
\begin{itemize}
\item equity payoff is $(y^i_h, y^i_h, y^i_\ell)$
\item the marketed subspace is contained in $\{w \in \R^3 \colon w_1 = w_2\}$
\item every security written on equity is constant across $s_1$ and $s_2$
\end{itemize}
\bigskip

If the entrepreneur deviates to low effort:
\begin{itemize}
\item the equity payoff becomes $(y^i_h, y^i_\ell, y^i_\ell)$
\item this vector is \emph{not} in the equilibrium marketed subspace ($w_1 \neq w_2$)
\end{itemize}
\bigskip

\alert{Partial spanning fails at $\overline{y}$}, and the RCPP price perception is ill-defined
\end{frame}

%====================================================================
\section{Constrained Pareto Efficiency}
%====================================================================

\begin{frame}\frametitle{The trade-off}
\textbf{The fundamental tension:}
\bigskip

Entrepreneurs who want to finance investment without large debt can issue equity:
\begin{itemize}
\item opens the way to risk-sharing and diversification
\item but means no longer receiving the full benefit of their effort
\item incentives to exert effort are diminished
\end{itemize}
\bigskip

\begin{block}{Question}
Do markets induce entrepreneurs to make the optimal trade-off between incentives and risk-sharing in their choice of debt and equity?
\end{block}
\end{frame}

%--------------------------------------------------------------------
\begin{frame}\frametitle{What does ``optimal'' mean?}
We compare the equilibrium to an allocation chosen by a planner facing the same frictions:
\bigskip

\begin{itemize}
\item the planner uses the same financial markets as the agents
\item the planner cannot dictate effort: entrepreneurs choose $e^i$ optimally given the financial structure imposed by the planner
\item the planner respects market clearing
\end{itemize}
\bigskip

This defines a notion of \emph{constrained} efficiency: better than no notion of efficiency, but weaker than the unconstrained first-best
\end{frame}

%--------------------------------------------------------------------
\begin{frame}\frametitle{Constrained feasible allocations}
\begin{definitionns}
$(\bm{c}, \bm{e}) = (c^i, e^i)_{i \in I}$ is \alert{constrained feasible} if there exist $(\bm{k}, \bm{z})$ such that:
\end{definitionns}

\begin{itemize}
\item[(i)] production: $y^i = f^i(k^i, e^i)$
\item[(ii)] optimal effort: $e^i \in \widetilde{e}^i(k^i, z^i, y^{-i})$
\item[(iii)] risk-sharing: $c^i_s = V_s(y) \cdot z^i$ for all $s \in S$
\item[(iv)] date-$0$ market clearing: $\sum_i (c^i_0 + k^i) = \sum_i \omega^i_0$
\item[(v)] portfolio market clearing:
\begin{equation*}
\sum_\ell z^\ell_j = 1 \text{ for $j \in I$}, \quad \sum_\ell z^\ell_j = 0 \text{ for $j \notin I$}
\end{equation*}
\end{itemize}
\end{frame}

%--------------------------------------------------------------------
\begin{frame}\frametitle{Constrained Pareto optimality}
\begin{definitionns}
A constrained feasible allocation $(\bm{c}, \bm{e})$ is \alert{constrained Pareto optimal} (CPO) if there is no other constrained feasible $(\widehat{\bm{c}}, \widehat{\bm{e}})$ with
\begin{equation*}
U^i(\widehat{c}^i, \widehat{e}^i) \geq U^i(c^i, e^i), \quad \text{for all $i \in I$}
\end{equation*}
and strict inequality for at least one $i$.
\end{definitionns}
\bigskip

\textbf{Question:} Is every RCPP equilibrium constrained Pareto optimal?
\end{frame}

%--------------------------------------------------------------------
\begin{frame}\frametitle{Why an RCPP equilibrium may fail to be CPO}
\textbf{Two potential sources of inefficiency:}
\bigskip

\begin{itemize}
\item \textbf{Externalities through effort.} Entrepreneur $i$'s effort affects the payoff of every security written on firm $i$'s output, hence the consumption of every holder. The equilibrium ignores this; a planner would internalize it.
\bigskip
\item \textbf{Nash interactions through outputs.} An RCPP equilibrium is implicitly a Nash equilibrium: each entrepreneur takes the others' decisions as given. A planner could co-ordinate on better effort levels.
\end{itemize}
\bigskip

The next section identifies a spanning condition that neutralizes both channels
\end{frame}

%====================================================================
\section{The Welfare Theorem}
%====================================================================

\begin{frame}\frametitle{No index-based securities}
\begin{assumption}
For every $i \in I_1$, $J^i_{-i} = \emptyset$: every security depends \emph{exclusively} on firm $i$'s output or \emph{not at all}
\end{assumption}
\bigskip

Under this assumption, agent $i$'s date-$1$ income decomposes:
\begin{equation*}
c^i_1 = m^i_1 + \sum_{j \in J^i} V_j(f^i(k^i, e^i)) z^i_j
\end{equation*}

where $m^i_1 := \sum_{j \in J_{-i}} V_j(y^{-i}) z^i_j$ is the \emph{outside income}
\bigskip

The other component is the \emph{inside income}, which depends on $i$'s effort
\end{frame}

%--------------------------------------------------------------------
\begin{frame}\frametitle{Outside subspace}
For every $i \in I$, define the outside operator $V_{-i}(y^{-i}) : \R^{J_{-i}} \to \R^S$ associated with the matrix $(V_j(y^{-i}))_{j \in J_{-i}}$
\bigskip

The \emph{outside subspace} is its range:
\begin{equation*}
\mathcal{V}_{-i}(y^{-i}) := \Range V_{-i}(y^{-i}) \subset \R^S
\end{equation*}
\bigskip

The choice of effort by entrepreneurs $\ell \neq i$ affects agent $i$ \emph{only through the outside subspace}
\bigskip

We need to limit how changes in others' efforts can affect this subspace
\end{frame}

%--------------------------------------------------------------------
\begin{frame}\frametitle{Strong partial spanning}
\begin{definitionns}[Strong partial spanning, SPS]
There is SPS at $\overline{y}$ if, for all alternative investments and efforts $(\bm{k}, \bm{e})$ with outputs $y^\ell = f^\ell(k^\ell, e^\ell)$:
\begin{equation*}
\mathcal{V}_{-i}(y^{-i}) \subset \mathcal{V}_{-i}(\overline{y}^{-i}), \quad \text{for all $i \in I_1$}
\end{equation*}
\end{definitionns}
\bigskip

The outside subspace, while it may change with others' efforts, stays within its equilibrium span
\bigskip

This rules out coordinated deviations creating new outside cash flows
\end{frame}

%--------------------------------------------------------------------
\begin{frame}\frametitle{SPS: comments}
\begin{itemize}
\item Even complete markets ($\mathcal{V}(\overline{y}) = \R^S$) do not automatically imply SPS, since $\mathcal{V}_{-i}(y^{-i}) \subsetneq \mathcal{V}(y)$
\bigskip
\item SPS holds if securities on the outputs of any subset of $I_1 \setminus \{i\}$ firms suffice to complete the markets
\bigskip
\item SPS holds in the two examples (single-factor; independent shocks with options): each firm spans its own subspace
\end{itemize}
\end{frame}

%--------------------------------------------------------------------
\begin{frame}\frametitle{The welfare theorem}
\begin{theoremns}[Constrained Pareto optimality of RCPP]
Let $\bigl\{ (\overline{c}^i, \overline{e}^i, \overline{k}^i, \overline{z}^i)_i, (\overline{y}^j)_j, \overline{q} \bigr\}$ be an RCPP equilibrium. Suppose:
\begin{itemize}
\item time-separable preferences
\item common knowledge of characteristics
\item no index-based securities
\item strong partial spanning at $\overline{y}$
\end{itemize}
Then $(\overline{c}^i, \overline{e}^i)_{i \in I}$ is constrained Pareto optimal.
\end{theoremns}
\end{frame}

%--------------------------------------------------------------------
\begin{frame}\frametitle{The welfare theorem: significance}
\textbf{The result is remarkable:}
\bigskip

Despite the moral-hazard friction, the decentralized equilibrium fully internalizes the externalities through the signalling channel
\bigskip

The entrepreneur, rationally anticipating that investors will read his financial decisions as signals of his effort, is led to choose precisely the effort level that the planner would have wanted
\bigskip

\textbf{Comparison with the previous part (Chapter on competitive price perceptions):}
\begin{itemize}
\item Same logic: Makowski-style price perceptions plus a spanning condition
\item Moral hazard changes the inferential content of financial choices, but not the welfare logic
\end{itemize}
\end{frame}

%====================================================================
\section{Sketch of Proof}
%====================================================================

\begin{frame}\frametitle{Proof: the strategy}
Suppose for contradiction that $(\bm{c}, \bm{e})$ is constrained feasible and Pareto-dominates $(\overline{\bm{c}}, \overline{\bm{e}})$
\bigskip

\textbf{Steps:}
\begin{enumerate}
\item Construct, for each agent $i$, an alternative portfolio $\widehat{z}^i$ that replicates the outside income $m^i_1$ at equilibrium outputs
\item Show that the optimal effort is invariant under this construction
\item Compute the implied date-$0$ consumption $\widehat{c}^i_0$
\item Use RCPP-optimality to conclude $c^i_0 \geq \widehat{c}^i_0$ for all $i$, with strict inequality for at least one
\item Aggregate to derive a contradiction with market clearing
\end{enumerate}
\end{frame}

%--------------------------------------------------------------------
\begin{frame}\frametitle{Proof: replicating portfolio}
For each $i \in I$, define $(\widehat{k}^i, \widehat{z}^i)$ by:
\begin{equation*}
\widehat{k}^i = k^i, \quad \widehat{z}^i_j = z^i_j \text{ for $j \in J^i$}
\end{equation*}
\bigskip

and $(\widehat{z}^i_j)_{j \in J_{-i}}$ such that
\begin{equation*}
\sum_{j \in J_{-i}} V_j(\overline{y}^{-i}) \widehat{z}^i_j = m^i_1
\end{equation*}
\bigskip

\textbf{Why is this feasible?} By the no-index assumption and SPS:
\begin{itemize}
\item the right-hand side $m^i_1$ belongs to $\mathcal{V}_{-i}(y^{-i})$
\item by SPS, $\mathcal{V}_{-i}(y^{-i}) \subset \mathcal{V}_{-i}(\overline{y}^{-i})$
\item so a portfolio of \emph{outside} securities exists that replicates $m^i_1$
\end{itemize}
\end{frame}

%--------------------------------------------------------------------
\begin{frame}\frametitle{Proof: effort invariance}
\begin{lemmans}[Effort invariance]
\begin{equation*}
\widetilde{e}^i(\widehat{k}^i, \widehat{z}^i, \overline{y}^{-i}) = \widetilde{e}^i(k^i, z^i, y^{-i})
\end{equation*}
\end{lemmans}
\bigskip

\textbf{Sketch.} For every $e^i$, the date-$1$ consumption under the alternative coincides with the original:
\begin{equation*}
\widehat{c}^i_1 = \sum_{j \in J^i} V_j(f^i(k^i, e^i)) z^i_j + m^i_1 = c^i_1
\end{equation*}
\bigskip

Since $u^i_1$ depends only on $(c^i_1, e^i)$, the argmax sets coincide
\end{frame}

%--------------------------------------------------------------------
\begin{frame}\frametitle{Proof: implied date-$0$ consumption}
By the RCPP price perception:
\begin{equation*}
\widehat{c}^i_0 = \omega^i_0 + \overline{\pi} \cdot f^i(k^i, \widehat{e}^i) - \overline{\pi} \cdot \widehat{c}^i_1 - k^i
\end{equation*}
\bigskip

Choosing $\widehat{e}^i = e^i$ (in the common argmax by Lemma) and using $\widehat{c}^i_1 = c^i_1$:
\begin{equation*}
\widehat{c}^i_0 = \omega^i_0 + \overline{\pi} \cdot y^i - \overline{\pi} \cdot c^i_1 - k^i
\end{equation*}
\end{frame}

%--------------------------------------------------------------------
\begin{frame}\frametitle{Proof: individual comparison}
The alternative $(\widehat{k}^i, \widehat{z}^i, \widehat{c}^i_0, c^i_1, e^i)$ is in the RCPP budget set
\bigskip

By RCPP-optimality of $(\overline{c}^i, \overline{e}^i)$:
\begin{equation*}
U^i(\overline{c}^i, \overline{e}^i) \geq U^i(\widehat{c}^i_0, c^i_1, e^i)
\end{equation*}
\bigskip

Combined with Pareto dominance:
\begin{equation*}
U^i(c^i_0, c^i_1, e^i) \geq U^i(\widehat{c}^i_0, c^i_1, e^i)
\end{equation*}
\bigskip

By time-separability and strict monotonicity of $u^i_0$:
\begin{equation*}
c^i_0 \geq \widehat{c}^i_0, \quad \text{for all $i \in I$}
\end{equation*}
with strict inequality for the agent $i^\ast$ enjoying strict utility improvement
\end{frame}

%--------------------------------------------------------------------
\begin{frame}\frametitle{Proof: aggregation and contradiction}
Summing over $i$:
\begin{equation*}
\sum_{i \in I} c^i_0 > \sum_{i \in I} \omega^i_0 + \overline{\pi} \cdot \sum_{i \in I} y^i - \overline{\pi} \cdot \sum_{i \in I} c^i_1 - \sum_{i \in I} k^i
\end{equation*}
\bigskip

By date-$0$ market clearing: $\sum_i (c^i_0 + k^i) = \sum_i \omega^i_0$
\bigskip

By portfolio market clearing: $\sum_i c^i_1 = \sum_i y^i$
\bigskip

Substituting yields $0 > 0$, a contradiction $\boite$
\end{frame}

%====================================================================
\section{Conclusion}
%====================================================================

\begin{frame}\frametitle{Conclusion}
\textbf{Constrained efficiency of the stock market} under moral hazard
\bigskip

Two behavioral assumptions:
\begin{enumerate}
\item \emph{Rationality of agents:} investors correctly infer the entrepreneur's effort from his observable financing decisions, and the entrepreneur is aware of this
\item \emph{Competitive markets:} an entrepreneur cannot influence the equilibrium state prices
\end{enumerate}
\bigskip

Plus structural conditions on the security structure (PS, SPS, no indices)
\end{frame}

%--------------------------------------------------------------------
\begin{frame}\frametitle{Intuition}
\begin{itemize}
\item The entrepreneur's asset holdings act as a \emph{signal} for his effort choice
\bigskip
\item Investors re-price the firm's securities based on their inference of effort
\bigskip
\item The entrepreneur, anticipating this, internalizes the inefficiency from lowered incentives through a lower price for his firm's stock
\bigskip
\item The externality at the heart of the moral-hazard problem is internalized through the price mechanism
\end{itemize}
\end{frame}

%--------------------------------------------------------------------
\begin{frame}\frametitle{Limitations}
\textbf{Calcagno--Wagner (2006):}
\begin{itemize}
\item The market allocation is constrained efficient \emph{only when} the entrepreneur has \emph{full initial ownership} of his firm
\item With dispersed initial ownership, the signalling mechanism breaks down
\end{itemize}
\bigskip

\textbf{Quiggin--Chambers (2006):}
\begin{itemize}
\item The market allocation is constrained efficient \emph{only when} the production technology has a \emph{single scalar input}
\item With multi-dimensional inputs, signalling cannot pin down the optimal composition
\end{itemize}
\end{frame}

%====================================================================
\section{Bibliography}
%====================================================================

\begin{frame}\frametitle{Main reference}
\begin{thebibliography}{8}
\beamertemplatearticlebibitems
    \bibitem{MagillQuinzii2002}
    Magill, M. and Quinzii, M.
    \newblock Capital market equilibrium with moral hazard.
    \newblock {\em Journal of Mathematical Economics} 38, 2002.
\end{thebibliography}
\bigskip

The exposition of this part is based on this paper
\end{frame}

%--------------------------------------------------------------------
\begin{frame}\frametitle{Earlier work and contemporary developments}
\begin{thebibliography}{8}
\begin{small}
\beamertemplatearticlebibitems
    \bibitem{KihlstromMatthews1990}
    Kihlstrom, R. E. and Matthews, S. A.
    \newblock Managerial incentives in an entrepreneurial stock market model.
    \newblock {\em Journal of Financial Intermediation} 1, 1990.
    \bibitem{Kocherlakota1998}
    Kocherlakota, N. R.
    \newblock The effects of moral hazard on asset prices when financial markets are complete.
    \newblock {\em Journal of Monetary Economics} 41, 1998.
    \bibitem{BisinGottardiRampini2008}
    Bisin, A., Gottardi, P. and Rampini, A.
    \newblock Managerial hedging and portfolio monitoring.
    \newblock {\em Journal of the European Economic Association} 6, 2008.
    \bibitem{AcharyaBisin2009}
    Acharya, V. V. and Bisin, A.
    \newblock Managerial hedging, equity ownership, and firm value.
    \newblock {\em RAND Journal of Economics} 40, 2009.
\end{small}
\end{thebibliography}
\end{frame}

%--------------------------------------------------------------------
\begin{frame}\frametitle{Comments and extensions on Magill--Quinzii (2002)}
\begin{thebibliography}{8}
\beamertemplatearticlebibitems
    \bibitem{CalcagnoWagner2006}
    Calcagno, R. and Wagner, W.
    \newblock Dispersed initial ownership and the efficiency of the stock market under moral hazard.
    \newblock {\em Journal of Mathematical Economics} 42, 2006.
    \bibitem{QuigginChambers2006}
    Quiggin, J. and Chambers, R.
    \newblock Capital market equilibrium with moral hazard and flexible technology.
    \newblock {\em Journal of Mathematical Economics} 42, 2006.
\end{thebibliography}
\end{frame}

\end{document}